% ===== PRÉAMBULE CANONIQUE — à coller VERBATIM en tête de CHAQUE .tex =====
\documentclass[11pt,a4paper]{article}
\usepackage[utf8]{inputenc}\usepackage[T1]{fontenc}\usepackage{lmodern}
\usepackage[french]{babel}
\usepackage{amsmath,amssymb,mathtools}\usepackage{icomma}
\usepackage{siunitx}\sisetup{locale=FR}  % \qty{}{}, \si{}, \ang{}
\usepackage{xcolor}\usepackage{colortbl}
\usepackage{tikz}\usepackage{pgfplots}\pgfplotsset{compat=1.18}
\usepackage[siunitx,europeanresistors]{circuitikz} % circuits (élec.)
\usepackage[most]{tcolorbox}\usepackage{enumitem}\usepackage{array,booktabs}
\usepackage[a4paper,margin=2.2cm]{geometry}
\usetikzlibrary{arrows.meta,calc,angles,quotes,positioning,patterns,%
  backgrounds,shapes.geometric,decorations.pathmorphing,decorations.markings,intersections}
\definecolor{primary}{RGB}{222,49,99}\definecolor{primarydark}{RGB}{150,22,55}
\definecolor{defcol}{RGB}{0,114,178}\definecolor{methcol}{RGB}{52,92,148}
\definecolor{excol}{RGB}{0,135,90}\definecolor{attncol}{RGB}{200,40,55}
\tcbset{boxbase/.style={enhanced,breakable,boxrule=0.9pt,arc=2.5pt,
  left=3mm,right=3mm,top=2.3mm,bottom=2.3mm,fonttitle=\bfseries,coltitle=white}}
% --- boîtes TUTORIEL (noms EXACTS, ne pas renommer) ---
\newtcolorbox{cle}[1][]{boxbase,colback=defcol!6,colframe=defcol,
  title={\textsc{À retenir}\ifx\relax#1\relax\relax\else\textmd{ : \itshape #1}\fi}}
\newtcolorbox{methode}[1][]{boxbase,colback=methcol!7,colframe=methcol,sharp corners,
  title={\textsc{Méthode}\ifx\relax#1\relax\relax\else\textmd{ --- \itshape #1}\fi}}
\newtcolorbox{exemplebox}[1][]{boxbase,colback=excol!5,colframe=excol,
  title={\textsc{Exemple}\ifx\relax#1\relax\relax\else\textmd{ : \itshape #1}\fi}}
\newtcolorbox{piege}[1][]{boxbase,colback=attncol!6,colframe=attncol,
  title={\textsc{Piège}\ifx\relax#1\relax\relax\else\textmd{ : \itshape #1}\fi}}
\newtcolorbox{remarque}[1][]{boxbase,colback=black!4,colframe=black!45,
  title={\textsc{Remarque}\ifx\relax#1\relax\relax\else\textmd{ : \itshape #1}\fi}}
% --- boîtes EXERCICE / CORRIGÉ (noms EXACTS, auto-comptées) ---
\newtcolorbox[auto counter]{exo}[1][]{boxbase,colback=primary!4,colframe=primary,
  title={\textsc{Exercice}~\thetcbcounter\ifx\relax#1\relax\relax\else\textmd{ --- \itshape #1}\fi}}
\newtcolorbox[auto counter]{corrige}[1][]{boxbase,colback=excol!4,colframe=excol,
  title={\textsc{Corrigé}~\thetcbcounter\ifx\relax#1\relax\relax\else\textmd{ --- \itshape #1}\fi}}
\newcommand{\vv}[1]{\vec{#1}}\newcommand{\ds}{\displaystyle}
\newcommand{\R}{\mathbb{R}}\newcommand{\N}{\mathbb{N}}\newcommand{\Z}{\mathbb{Z}}
% (un \usepackage{fancyhdr}+en-tête est possible ; il est ignoré côté web)
% =========================================================================
\begin{document}
\sisetup{per-mode=symbol}

\begin{exo}[la vraie vitesse moyenne]
Un automobiliste parcourt \qty{120}{\kilo\metre} à \qty{40}{\kilo\metre\per\hour}, puis
\qty{120}{\kilo\metre} à \qty{60}{\kilo\metre\per\hour}.
\begin{enumerate}[label=\alph*)]
  \item Calcule la durée de chaque tronçon.
  \item En déduis sa vitesse moyenne sur l'ensemble du trajet.
  \item Pourquoi la réponse n'est-elle pas \qty{50}{\kilo\metre\per\hour} ?
\end{enumerate}
\end{exo}

\begin{exo}[dérivée d'un mouvement cubique]
La position d'un mobile est $x(t)=t^3-6t^2+9t$ (en \si{\metre}, $t\ge0$ en \si{\second}).
\begin{enumerate}[label=\alph*)]
  \item Donne $v(t)$.
  \item À quels instants la vitesse s'annule-t-elle ?
  \item Donne $a(t)$ et calcule sa valeur à ces deux instants.
\end{enumerate}
\end{exo}

\begin{exo}[remonter par intégration]
Un mobile a pour vitesse $v(t)=6t+2$ (\si{\metre\per\second}) et part de $x_0=\qty{5}{\metre}$ à $t=0$.
\begin{enumerate}[label=\alph*)]
  \item Détermine l'équation horaire $x(t)$.
  \item Où se trouve-t-il à $t=\qty{4}{\second}$ ?
  \item Quelle est son accélération ?
\end{enumerate}
\end{exo}

\begin{exo}[lecture d'un graphe de vitesse]
Le graphe donne la vitesse d'un véhicule sur $Ox$.
\begin{center}
\begin{tikzpicture}
\begin{axis}[width=9cm,height=4.6cm,axis lines=middle,xlabel={$t\ (\si{\second})$},ylabel={$v\ (\si{\metre\per\second})$},
  xmin=0,xmax=7.6,ymin=0,ymax=9.5,xtick={0,1,2,3,4,5,6,7},ytick={0,2,4,6,8},grid=both,
  grid style={black!12},clip=false]
  \addplot[primary,line width=1.3pt] coordinates {(0,0)(2,8)(5,8)(7,0)};
\end{axis}
\end{tikzpicture}
\end{center}
\begin{enumerate}[label=\alph*)]
  \item Calcule l'accélération sur chacune des trois phases $[0;2]$, $[2;5]$, $[5;7]$~\si{\second}.
  \item En déduis la distance totale parcourue (aire sous la courbe).
\end{enumerate}
\end{exo}

\begin{exo}[intégrer une vitesse non affine]
Un mobile parti de l'origine a pour vitesse $v(t)=3t^2$ (\si{\metre\per\second}).
\begin{enumerate}[label=\alph*)]
  \item Quelle distance parcourt-il entre $t=0$ et $t=\qty{2}{\second}$ ?
  \item Donne son accélération $a(t)$ : est-elle constante ?
\end{enumerate}
\end{exo}

\end{document}
